% !TeX root = ../../Skript.tex
\cohead{\Large\textbf{Hyperbeln}}
\fakesubsection{Hyperbeln}

\begin{Exercise}[title={}, label=HyperbelA1]
	
Beschreibe die Form des Schaubildes von Funktionen vom Typ \(f(x)=x^{-n},\ n\in\N\) 

Wie transformieren die Parameter \( a \), \( b \) und \( c \) das Schaubild: \( f(x) = a \cdot (x - b)^{-n} + c \)?  
\end{Exercise}

\begin{Exercise}[title={}, label=HyperbelA2]
	
Gegeben ist \( f(x) = 2 \cdot (x - 1)^{-2} + 3 \). Berechne die Funktionswerte für:  
\begin{enumerate}[label=\alph*)]
	\item \( x_1 = 2 \) 
	\item \( x_2 = 0 \)  
	\item \( x_3 = 1.5 \)
\end{enumerate}
\end{Exercise}  

\begin{Exercise}[title={}, label=HyperbelA3]
	
Gegeben ist die Funktion \( f(x) = -4 \cdot (x - 3)^{-1} + 2 \).  
\begin{enumerate}[label=\alph*)]
	\item Bestimme die waagrechte und die horizontale Asymptote.
	\item Erläutere, was diese Asymptoten über das Verhalten der Funktion aussagen. 
\end{enumerate}
\end{Exercise}

\begin{Exercise}[title={}, label=HyperbelA4]
	
Gegeben ist die Funktion \( f(x) = a \cdot (x - 2)^{-3} + 5 \).  
\begin{enumerate}[label=\alph*)]
	\item Beschreibe, wie sich der Graph verändert, wenn \( a \) verdoppelt wird.  
	\item Was passiert, wenn \( c=5 \) auf \( -2 \) geändert wird? 
	\item Welchen Einfluss hat \( n \) auf die Form der Funktion?
\end{enumerate}
\end{Exercise}  

\begin{Exercise}[title={}, label=HyperbelA5] 
	
Untersuche die Funktion \( f(x) = \frac{3}{(x + 1)^2} - 2 \). 
\begin{enumerate}[label=\alph*)]
	\item Hat die Funktion Nullstellen? Begründe. 
	\item Wo hat die Funktion Definitionslücken?
\end{enumerate} 
\end{Exercise}  

\begin{Exercise}[title={}, label=HyperbelA6]
	
Eine Funktion des Typs \( f(x) = a \cdot (x - b)^{-2} + c \) hat folgende Eigenschaften: 
\begin{itemize}
	\item Eine vertikale Asymptote bei \( x = 4 \).
	\item Eine waagrechte Asymptote bei \( y = -1 \).
	\item Der Funktionswert an der Stelle \( x = 5 \) beträgt \( 3 \).
\end{itemize} 
Bestimme \( a \), \( b \) und \( c \).
\end{Exercise}

\begin{Exercise}[title={}, label=HyperbelAllgemein]
	
	Gib die Definitions- und Wertemenge an, gib die Asymptoten an und skizziere das Schaubild.
	
	\begin{minipage}{\textwidth}
		\begin{minipage}{0.5\textwidth}
			\begin{enumerate}[label=\alph*)]
				\item \(f(x)=x^{-2}+1\)
				\item \(f(x)=\left(x+3\right)^{-4}\)
				\item \(f(x)=\frac{3}{x}\)
				\item \(f(x)=-x^{-1}\)
				\item \(f(x)=\frac{2}{x-4}-2\)
				\item \(f(x)=-0,5\left(x+1\right)^{-1}+2\)
			\end{enumerate}
		\end{minipage}%
		\begin{minipage}{0.5\textwidth}
			\begin{enumerate}[label=\alph*)]
				\setcounter{enumi}{6}
				\item \(f(x)=-0,5\left(x+1,5\right)^{-3}\)
				\item \(f(x)=-2+\frac{1}{4x^2}\)
				\item \(f(x)=5-\left(x+2\right)^{-2}\)
				\item \(f(x)=\frac{\left(x-2\right)^{-1}+4}{2}\)
				\item \(f(x)=\frac{1}{5}\left(x-\frac{3}{2}\right)^{-1}+\frac{7}{5}\)
				\item \(f(x)=-2-3\left(x+4\right)^{-5}\)
			\end{enumerate}
		\end{minipage}%
	\end{minipage}
\end{Exercise}
\newpage
%%%%%%%%%%%%%%%%%%%%%%%%%%%%%%%%%%%%%%%%%
\begin{Answer}[ref=HyperbelA1]
	Schaubilder vom Typ \(f(x)=x^{-n}\) verlaufen durch den Punkt \(P(1\vert 1)\). Für größere Werte von \(x\) nähert sich das Schaubild immer mehr der \(x\)-Achse an ohne diese jemals zu erreichen. Für \(x\) gegen Null gehen die Funktionswerte gegen Unendlich, so dass sich das Schaubild der \(y\)-Achse beliebig nahe nähert. \(x=0\) ist eine Definitionslücke, da die Division durch Null nicht definiert ist. Für positive \(x\)-Werte fällt das Schaubild also von Unendlich kommend immer mehr gegen Null. 
	
	Ist \(n\) ungerade, so erhält man das Schaubild für negative \(x\) durch eine Punktspiegelung am Ursprung. Ist \(n\) gerade, so erhält man das Schaubild für negative \(x\) durch eine Achsenspiegelung an der \(y\)-Achse.
\end{Answer}
\begin{Answer}[ref=HyperbelA2]
	\begin{enumerate}[label=\alph*)]
		\item \(f(2)=2(2-1)^{-2}+3=2+3=5\)
		\item \(f(0)=2(0-1)^{-2}+3=2+3=5\)
		\item \(f(1,5)=2(1,5-1)^{-2}+3=2\cdot0,5^{-2}+3=8+3=11\)
	\end{enumerate}
\end{Answer}
\begin{Answer}[ref=HyperbelA3]
	\begin{enumerate}[label=\alph*)]
		\item waagrechte Asymptote: \(y=2\), senkrechte Asymptote \(x=3\)
		\item Für \(x\to\pm\infty\) nähert sich das Schaubild der waagrechten Asymptoten an und für \(x\to 3\) gehen die Funktionswerte gegen \(\pm\infty\).
	\end{enumerate}
\end{Answer}
\begin{Answer}[ref=HyperbelA4]
	\begin{enumerate}[label=\alph*)]
		\item \(a\) ist der Streckfaktor (relativ zur waagrechten Asymptoten), d.h. wenn man \(a\) verdoppelt, verdoppelt sich der Abstand jedes Punktes des Schaubildes zur waagrechten Asymptoten \(y=5\).
		\item Das Schaubild wird um 7 Einheiten nach unten verschoben.
		\item Für gerade \(n\) ist das Schaubild achsensymmetrisch zur senkrechten Asymptoten, für ungerade \(n\) punktsymmetrisch zum Schnittpunkt von waagrechter und senkrechter Asymptote.
	\end{enumerate}

\end{Answer}
\begin{Answer}[ref=HyperbelA5]
	\begin{enumerate}[label=\alph*)]
		\item Ja, die Funktion hat 2 Nullstellen, da das Schaubild für \(x\to \infty\) gegen die waagrechte Asymptote \(y=-2\) geht \(\left(\lim\limits_{x\to\infty}=-2\right)\) und für \(x\to -1\) gilt \(f(x)\to\infty\), da der Streckfaktor \(a=3\) positiv ist. Es muss also eine Nullstelle zwischen rechts von \(x=-1\) geben. Da \(n\) gerade ist, ist das Schaubild achsensymmetrisch zu \(x=-1\), wodurch sich die 2. Nullstelle ergibt. 
		\item Die Funktion hat eine Definitionslücke bei \(x=-1\), (da die Division durch Null nicht definiert ist.)
	\end{enumerate}

\end{Answer}
\begin{Answer}[ref=HyperbelA6]
	\begin{itemize}
		\item Vertikale Asymptote bei \(x=4\), also \(b=4\).
		\item Waagrechte Asymptote bei \(y=-1\), also \(c=-1\).
		\item Aus \(f(5)=3\) ergibt sich \(a\left(5-4\right)^{-2}-1=3\) und damit \(a=4\)
	\end{itemize}
	\( f(x) = 4 \cdot (x - 4)^{-2} -1 \)
\end{Answer}
\begin{Answer}[ref=HyperbelAllgemein]
	
	\begin{minipage}{\textwidth}
		\begin{minipage}{0.5\textwidth}
			\begin{enumerate}[label=\alph*)]
				\item \(f(x)=x^{-2}+1\)
				
				\(D_f=\R\setminus\{0\},\ W_f=]1;\infty[=\{x\in\R\vert x>1\}\)
				
				Asymptoten \(y=1\) und \(x=0\)
				
				\adjustbox{valign=t, padding=0ex 0ex 0ex 0ex}{\begin{minipage}{\linewidth}
						\includegraphics[width=6cm]{\ganzFkt/pics/HyperbelA1picA.png}
				\end{minipage}}%			
				
				\item \(f(x)=\left(x+3\right)^{-4}\)
				
				\(D_f=\R\setminus\{-3\},\ W_f=]0;\infty[=\R^+\)
				
				Asymptoten \(y=0\) und \(x=-3\)
				
				\adjustbox{valign=t, padding=0ex 0ex 0ex 0ex}{\begin{minipage}{\linewidth}
						\includegraphics[width=6cm]{\ganzFkt/pics/HyperbelA1picB.png}
				\end{minipage}}%
				\item \(f(x)=\frac{3}{x}\)
				
				\(D_f=\R\setminus\{0\},\ W_f=\R\setminus\{0\}\)
				
				Asymptoten \(y=0\) und \(x=0\)
				
				\adjustbox{valign=t, padding=0ex 0ex 0ex 0ex}{\begin{minipage}{\linewidth}
						\includegraphics[width=6cm]{\ganzFkt/pics/HyperbelA1picC.png}
				\end{minipage}}%
			\end{enumerate}
		\end{minipage}%
		\begin{minipage}{0.5\textwidth}
			\begin{enumerate}[label=\alph*)]
				\setcounter{enumi}{3}
				\item \(f(x)=-x^{-1}\)
				
				\(D_f=\R\setminus\{0\},\ W_f=\R\setminus\{0\}\)
				
				Asymptoten \(y=0\) und \(x=0\)
				
				\adjustbox{valign=t, padding=0ex 0ex 0ex 0ex}{\begin{minipage}{\linewidth}
						\includegraphics[width=6cm]{\ganzFkt/pics/HyperbelA1picD.png}
				\end{minipage}}%
				\item \(f(x)=\frac{2}{x-4}-2\)
				
				\(D_f=\R\setminus\{4\},\ W_f=\R\setminus\{-2\}\)
				
				Asymptoten \(y=-2\) und \(x=4\)
				
				\adjustbox{valign=t, padding=0ex 0ex 0ex 0ex}{\begin{minipage}{\linewidth}
						\includegraphics[width=6cm]{\ganzFkt/pics/HyperbelA1picE.png}
				\end{minipage}}%
				\item \(f(x)=-0,5\left(x+1\right)^{-1}+2\)
				
				\(D_f=\R\setminus\{-1\},\ W_f=\R\setminus\{2\}\)
				
				Asymptoten \(y=2\) und \(x=-1\)
				
				\adjustbox{valign=t, padding=0ex 0ex 0ex 0ex}{\begin{minipage}{\linewidth}
						\includegraphics[width=6cm]{\ganzFkt/pics/HyperbelA1picF.png}
				\end{minipage}}%
			\end{enumerate}
		\end{minipage}%
	\end{minipage}
	
		
	\begin{minipage}{\textwidth}
		\begin{minipage}{0.5\textwidth}
			\begin{enumerate}[label=\alph*)]
				\setcounter{enumi}{6}
				\item \(f(x)=-0,5\left(x+1,5\right)^{-3}\)
				
				\(D_f=\R\setminus\{-1,5\},\ W_f=\R\setminus\{0\}\)
				
				Asymptoten \(y=0\) und \(x=-1,5\)
				
				\adjustbox{valign=t, padding=0ex 0ex 0ex 0ex}{\begin{minipage}{\linewidth}
						\includegraphics[width=6cm]{\ganzFkt/pics/HyperbelA1picG.png}
				\end{minipage}}%
				\item \(f(x)=-2+\frac{1}{4x^2}\)
				
				\(D_f=\R\setminus\{0\},\ W_f=\{x\in\R\vert x>-2\}\)
				
				Asymptoten \(y=-2\) und \(x=0\)
				
				\adjustbox{valign=t, padding=0ex 0ex 0ex 0ex}{\begin{minipage}{\linewidth}
						\includegraphics[width=6cm]{\ganzFkt/pics/HyperbelA1picH.png}
				\end{minipage}}%
				\item \(f(x)=5-\left(x+2\right)^{-2}\)
				
				\(D_f=\R\setminus\{-2\},\ W_f=]-\infty;5[\)
				
				Asymptoten \(y=5\) und \(x=-2\)
				
				\adjustbox{valign=t, padding=0ex 0ex 0ex 0ex}{\begin{minipage}{\linewidth}
						\includegraphics[width=6cm]{\ganzFkt/pics/HyperbelA1picI.png}
				\end{minipage}}%

			\end{enumerate}
		\end{minipage}%
		\begin{minipage}{0.5\textwidth}
			\begin{enumerate}[label=\alph*)]
				\setcounter{enumi}{9}
				\item \(f(x)=\frac{\left(x-2\right)^{-1}+4}{2}\)
				
				\(D_f=\R\setminus\{2\},\ W_f=]2;\infty[=\{x\in\R\vert x>2\}\)
				
				Asymptoten \(y=2\) und \(x=2\)
				
				\adjustbox{valign=t, padding=0ex 0ex 0ex 0ex}{\begin{minipage}{\linewidth}
						\includegraphics[width=6cm]{\ganzFkt/pics/HyperbelA1picJ.png}
				\end{minipage}}%
				\item \(f(x)=\frac{1}{5}\left(x-\frac{3}{2}\right)^{-1}+\frac{7}{5}\)
				
				\(D_f=\R\setminus\{\frac{3}{2}\},\ W_f=\R\setminus\{\frac{7}{5}\}\)
				
				Asymptoten \(y=\frac{7}{5}\) und \(x=\frac{3}{2}\)
				
				\adjustbox{valign=t, padding=0ex 0ex 0ex 0ex}{\begin{minipage}{\linewidth}
						\includegraphics[width=6cm]{\ganzFkt/pics/HyperbelA1picK.png}
				\end{minipage}}%
				\item \(f(x)=-2-3\left(x+4\right)^{-5}\)
				
				\(D_f=\R\setminus\{-4\},\ W_f=\R\setminus\{2\}\)
				
				Asymptoten \(y=2\) und \(x=-4\)
				
				\adjustbox{valign=t, padding=0ex 0ex 0ex 0ex}{\begin{minipage}{\linewidth}
						\includegraphics[width=6cm]{\ganzFkt/pics/HyperbelA1picL.png}
				\end{minipage}}%
			\end{enumerate}
		\end{minipage}%
	\end{minipage}
\end{Answer}